\\documentclass[12pt,a4paper]{article}
\\usepackage[utf8]{inputenc}
\\usepackage[vietnamese]{babel}
\\usepackage{amsmath}
\\usepackage{amsfonts}
\\usepackage{amssymb}
\\usepackage{graphicx}
\\usepackage{booktabs}
\\usepackage{hyperref}
\\usepackage{geometry}
\\usepackage{titlesec}
\\usepackage{enumitem}

\\geometry{margin=1in}

\\titleformat{\\section}{\\normalfont\\Large\\bfseries\\color{blue}}{\\thesection}{1em}{}
\\titleformat{\\subsection}{\\normalfont\\large\\bfseries\\color{darkgray}}{\\thesubsection}{1em}{}

\\title{\\textbf{Kế hoạch & Lộ trình Thực hiện Dự án (1 Tuần)\\\\Xây dựng IDS với Trí tuệ Nhân tạo}\\\\Nhóm 10 - Bộ môn An toàn Thông tin}
\\author{Đại học Quốc gia TP.HCM}
\\date{Tháng 6, Năm 2026}

\\begin{document}

\\maketitle

\\section{Phân chia Vai trò thành viên (5 Thành viên)}
Dự án được phân rã thành các mảng công việc chuyên biệt phù hợp với năng lực của từng thành viên:

\\begin{itemize}
    \\item \\textbf{Thành viên 1 (Nhóm trưởng / AI Engineer):}
    \\begin{itemize}
        \\item Chịu trách nhiệm điều phối chung, tích hợp mã nguồn của cả nhóm.
        \\item Nghiên cứu cấu hình mô hình học máy: Random Forest và XGBoost (`src/models.py`).
        \\item Tinh chỉnh siêu tham số (hyperparameter tuning) để đạt độ chính xác cao nhất.
    \\end{itemize}
    
    \\item \\textbf{Thành viên 2 (Data Engineer / Analyst):}
    \\begin{itemize}
        \\item Phụ trách thu thập tập dữ liệu từ Kaggle (CICIDS2017/CSE-CIC-IDS2018).
        \\item Thực hiện phân tích dữ liệu khám phá (EDA), trực quan hóa phân phối nhãn và tương quan thuộc tính.
        \\item Viết mã tiền xử lý, xử lý giá trị rỗng/vô cùng và cân bằng dữ liệu (`src/preprocessing.py`).
    \\end{itemize}
    
    \\item \\textbf{Thành viên 3 (Network Security Specialist):}
    \\begin{itemize}
        \\item Nghiên cứu các đặc trưng mạng (Network flow features).
        \\item Thu thập tập dữ liệu kiểm thử độc lập từ nguồn ngoài (ví dụ: UNSW-NB15) để chạy kiểm thử chéo (`src/evaluate.py`).
        \\item Cấu hình Npcap và chạy thực nghiệm sniffer trực tiếp (`src/live_sniffer.py`) để phát hiện tấn công.
    \\end{itemize}
    
    \\item \\textbf{Thành viên 4 (Frontend / Dashboard Developer):}
    \\begin{itemize}
        \\item Thiết kế và lập trình giao diện Web bằng Streamlit (`src/dashboard.py`).
        \\item Tích hợp biểu đồ thống kê kết quả phân tích file lưu lượng và tạo module mô phỏng log xâm nhập realtime.
    \\end{itemize}
    
    \\item \\textbf{Thành viên 5 (Technical Writer / Presenter):}
    \\begin{itemize}
        \\item Tổng hợp tài liệu, viết báo cáo học thuật bằng LaTeX hoặc Word.
        \\item Chuẩn bị slide thuyết trình của nhóm, căn chỉnh bố cục chuyên nghiệp.
        \\item Tổ chức các buổi chạy thử nghiệm thuyết trình thử (mock defense) trước ngày bảo vệ.
    \\end{itemize}
\\end{itemize}

\\newpage
\\section{Lộ trình Thực hiện Chi tiết (Roadmap 7 Ngày)}
Dưới đây là kế hoạch chi tiết từng ngày phục vụ buổi báo cáo sau 1 tuần:

\\subsection*{Ngày 1: Thiết lập Môi trường và EDA dữ liệu}
\\begin{itemize}
    \\item \\textbf{Công việc:} 
    \\begin{itemize}
        \\item Cả nhóm chạy file `run_project.py` (Lựa chọn 1) để cài đặt đồng bộ các thư viện phụ thuộc.
        \\item Thành viên 2 tải dataset mẫu đã làm sạch từ Kaggle và lưu vào `data/raw/`.
        \\item Thành viên 5 dựng cấu trúc tài liệu báo cáo và slide thuyết trình.
    \\end{itemize}
    \\item \\textbf{Kết quả bàn giao:} Môi trường chạy thành công, có file dữ liệu thật tại máy của các thành viên.
\\end{itemize}

\\subsection*{Ngày 2: Tiền xử lý dữ liệu và Lựa chọn thuộc tính}
\\begin{itemize}
    \\item \\textbf{Công việc:}
    \\begin{itemize}
        \\item Thành viên 2 thực hiện làm sạch dữ liệu, chuẩn hóa scaler và lưu lại file `scaler.pkl`.
        \\item Thành viên 1 viết phân tích lựa chọn 15 thuộc tính mạng cốt lõi giúp hệ thống chạy nhanh.
        \\item Thành viên 5 cập nhật chương lý thuyết về xử lý mất cân bằng dữ liệu (Class Imbalance).
    \\end{itemize}
    \\item \\textbf{Kết quả bàn giao:} Hoàn thiện module `preprocessing.py` hoạt động ổn định trên tệp lớn.
\\end{itemize}

\\subsection*{Ngày 3: Huấn luyện và So sánh các thuật toán AI}
\\begin{itemize}
    \\item \\textbf{Công việc:}
    \\begin{itemize}
        \\item Thành viên 1 chạy `train.py` để huấn luyện Random Forest, XGBoost và Isolation Forest.
        \\item Trích xuất các chỉ số Accuracy, Precision, Recall, F1-Score và vẽ Confusion Matrix của các mô hình.
        \\item Thành viên 5 lấy các kết quả bảng biểu này để đưa vào chương kết quả nghiên cứu trong báo cáo.
    \\end{itemize}
    \\item \\textbf{Kết quả bàn giao:} Các tệp mô hình đã lưu `.pkl` trong thư mục `models/` sẵn sàng cho việc sử dụng.
\\end{itemize}

\\subsection*{Ngày 4: Xây dựng Dashboard giao diện Web tương tác}
\\begin{itemize}
    \\item \\textbf{Công việc:}
    \\begin{itemize}
        \\item Thành viên 4 lập trình Web UI bằng Streamlit (`dashboard.py`).
        \\item Tích hợp tính năng upload file CSV và hiển thị biểu đồ tròn phân tích lưu lượng.
        \\item Thành viên 1 hỗ trợ liên kết mô hình AI đã train từ Ngày 3 vào dashboard để chạy dự đoán trực tiếp trên Web.
    \\end{itemize}
    \\item \\textbf{Kết quả bàn giao:} Giao diện Web Streamlit khởi chạy thành công qua lệnh `streamlit run src/dashboard.py`.
\\end{itemize}

\\subsection*{Ngày 5: Thực nghiệm Kiểm thử chéo và Sniffer Realtime}
\\begin{itemize}
    \\item \\textbf{Công việc:}
    \\begin{itemize}
        \\item Thành viên 3 tải dữ liệu test từ UNSW-NB15 dạng CSV và chạy `evaluate.py` để kiểm thử độ tổng quát.
        \\item Thành viên 3 cấu hình card mạng và chạy thử `live_sniffer.py`, thực hiện quét cổng (nmap) giả lập để kiểm tra phản hồi của hệ thống.
        \\item Thành viên 5 chụp ảnh/quay video màn hình quá trình cảnh báo realtime làm tài liệu thuyết minh.
    \\end{itemize}
    \\item \\textbf{Kết quả bàn giao:} Video/Ảnh chụp demo thực tế chạy sniffer và bảng số liệu đánh giá chéo.
\\end{itemize}

\\subsection*{Ngày 6: Hoàn thiện Báo cáo khoa học và Slide thuyết trình}
\\begin{itemize}
    \\item \\textbf{Công việc:}
    \\begin{itemize}
        \\item Thành viên 5 hoàn thiện tài liệu báo cáo đầy đủ các chương (Mở đầu, Cơ sở lý thuyết, Thực nghiệm, Kết quả, Kết luận).
        \\item Cả nhóm cùng rà soát lỗi chính tả và các phần lập luận kỹ thuật.
        \\item Đóng gói mã nguồn sạch sẽ, loại bỏ file log rác.
    \\end{itemize}
    \\item \\textbf{Kết quả bàn giao:} File báo cáo dạng PDF (compile từ LaTeX) và slide PowerPoint chuyên nghiệp.
\\end{itemize}

\\subsection*{Ngày 7: Tổng duyệt Demo và Thuyết trình thử}
\\begin{itemize}
    \\item \\textbf{Công việc:}
    \\begin{itemize}
        \\item Cả nhóm tổ chức họp trực tuyến/trực tiếp chạy thử Demo Dashboard Streamlit.
        \\item Thành viên 5 chạy slide và thuyết trình thử, các thành viên khác đặt câu hỏi phản biện giả lập.
        \\item Nộp bài tập lớn lên hệ thống quản lý học tập.
    \\end{itemize}
    \\item \\textbf{Kết quả bàn giao:} Dự án hoàn thiện 100%, sẵn sàng cho buổi bảo vệ chính thức trên lớp.
\\end{itemize}

\\end{document}
